\documentclass{jsarticle}
\usepackage{amsmath,amssymb}

\begin{document}

\begin{center}
\Large \bf
電子計算機及び実習２　課題1
\end{center}

\begin{flushright}
創域理工学部　数理科学科　6122111　三森　誠真
\end{flushright}

\bigskip
\bigskip

\textbf{Exercise 4.1}

(1) $p \equiv 1 \pmod{1000} を満たす最小の素数 p と 2 番目に小さな素数 q の積 n = pq$ を求めよ.

\bigskip

$1 \equiv 1 \pmod{1000}, \ 1001 \equiv 1 \pmod{1000}, \ 2001 \equiv 1 \pmod{1000}, \ 3001 \equiv 1 \pmod{1000},$

$4001 \equiv 1 \pmod{1000}, \dotsb\dotsb$ であり、

$1, \ 1001 = 7\cdot11\cdot13, \ 2001 = 3\cdot23\cdot29$は素数ではなく、$3001, \ 4001$は素数なので、

$p = 3001, \ q =4001$である。また、$n = pq =12007001$

\bigskip
\bigskip

(2) $上の n と e = 1001 に対して, \ ed ≡ 1 \pmod{\varphi(n)} となる最小の正整数 d $を求めよ.

\bigskip

まず、$\varphi(n) = (p-1)(q-1) = 12000000なので、1001d'≡1 \pmod{12000000}となるような最小の自然数d'$を求めればよい。上の式は、ある自然数$cを用いて1001d' - 12000000c = 1$ -------- \textcircled{\scriptsize 1} と表せるから、この不定方程式の解$d'$をユークリッドの互除法を用いて求めていく。

\bigskip

 \ \ \ \ \ \ \ \ \ \ \ \ \ \ \ \ \ \ \ $12000000 = 1001\cdot11988 + 12$

 \ \ \ \ \ \ \ \ \ \ \ \ \ \ \ \ \ \ \ \ \ \ \ \ \ $1001 = 12\cdot83 + 5$

 \ \ \ \ \ \ \ \ \ \ \ \ \ \ \ \ \ \ \ \ \ \ \ \ \ \ \ \ $12 = 5\cdot2 + 2$

 \ \ \ \ \ \ \ \ \ \ \ \ \ \ \ \ \ \ \ \ \ \ \ \ \ \ \ \ \ $5 = 2\cdot2 + 1$ \ \ \ \ \ \ \ \ \ \ なので、

\begin{align*}
1
&=5-2(12-5\cdot2)\\
&=5\cdot5-2\cdot12\\
&=5(1001-12\cdot83)-2\cdot12\\
&=1001\cdot5-12\cdot417\\
&=1001\cdot5-417(12000000-1001\cdot11988)\\
&=1001\cdot4999001+12000000\cdot(-417) ────── \ \textcircled{\scriptsize 2}
\end{align*}

\textcircled{\scriptsize 1}, \textcircled{\scriptsize 2}より $1001 (d'-4999001) = 12000000 (c-417)$

$1001と12000000は互いに素なので、ある整数kを用いて d'-4999001=12000000k 即ち d'=12000000k + 4999001と表せる。よって、kが整数であることから、最小の自然数d'を考えるとd=4999001$

\bigskip
\bigskip

(3) $上の (n, e) を暗号化伴, (n, d) を復号伴とした際, 平文 M_1 = 2023 の暗号化した結果 C_1$
を求めよ. 

 \ \ \ \ \ $また, 暗号文 C_2 = 5716652 を復号した結果 M_2$ を求めよ.

\bigskip

(計算にはWolfram Alphaを用いた)

$C_1 = {M_1}^e \pmod{n} = (2023)^{1001} \pmod{12007001} = 1502353$

$M_2 = {C_2}^d \pmod{n} = (5716652)^{4999001} \pmod{12007001} = 777$



\end{document}
